\begin{frame}{Naviguer dans un fichier XML}
    \begin{itemize}
        \item La hiérarchie d'un document XML se présente sous forme d'arbre constitué de \textcolor{red}{noeuds}
        \item Le langage XPath permet de sélectionner un noeud ou un ensemble de noeuds.
    \end{itemize}
    \textbf{Objectif:} Atteindre des fragments du document ayant certaines propriétés spécifiées dans une expression XPath
\end{frame}

\begin{frame}{Naviguer dans un fichier XML}
    \begin{figure}
        \centering
        \includegraphics[width=\linewidth]{03 - XPath/arbre_xml.png}
    \end{figure}
\end{frame}

\begin{frame}{Relations "généalogiques" entre les noeuds}
    \begin{itemize}
        \item \textbf{Parent =} Noeud \textcolor{red}{directement} au-dessus.
        \item \textbf{Enfant =} Noeud \textcolor{red}{directement} en dessous dans la hiérarchie.
        \item \textbf{Descendant =} Noeuds situés en dessous.
        \item \textbf{Ancêtre =} Noeuds situés au dessus.
        \item \textbf{Frère = } Noeuds situés au même niveau.
    \end{itemize}
\end{frame}

\begin{frame}{Notion central: chemin}
    \begin{itemize}
        \item Une expression (ou requête) décrit le chemin emprunté dans l'arbre pour atteindre les données qui nous intéressent.
        \item Un chemin peut être:
        \begin{itemize}
            \item \textcolor{red}{Absolu =} à partir de la racine.
            \item \textcolor{red}{Relatif =} à partir de l'endroit où l'on est dans le chemin (= le noeud courant).
        \end{itemize}
    \end{itemize}
\end{frame}

\begin{frame}{Notion central: chemin}
    \begin{block}{Chemin absolu}
        \begin{itemize}
            \item Commence par "/" suivie du noeud racine XML
            \item Connaissance exacte de la structure de l'arbre manipulé
            \item Chemin complet conduisant aux données à atteindre
        \end{itemize}
    \end{block}
    \begin{block}{Exemple: récupérer le titre du poème}
        \begin{itemize}
            \item Étape 1: noeud \textbf{TEI}
            \item Étape 2: noeud enfant \textbf{teiHeader}
            \item Étape 3: noeud enfant \textbf{fileDesc}
            \item Étape 4: noeud enfant \textbf{titleStmt}
            \item Étape 5: noeud enfant \textbf{title}
        \end{itemize}
        \textcolor{blue}{\texttt{/Tei/teiHeader/fileDesc/titleStmt/title}}
    \end{block}
\end{frame}

\begin{frame}{Notion central: chemin}
    \begin{block}{Chemin relatif}
        \begin{itemize}
            \item Le noeud de départ peut être n'importe quel noeud de l'arbre
            \item Commence par "//"
        \end{itemize}
    \end{block}
    \begin{block}{Exemple: récupérer le titre du poème}
        \vspace{0.5em}
        \textcolor{blue}{\texttt{//fileDesc/title}}
    \end{block}
\end{frame}
